%%%%%%%%%%%%%%%%%%%%%%%%%%%%%%%%%%%%%%%%%%%%%%%%%%%%%%%%%%%%
% 10.2 LINEAR PROGRAMMING
% The Composition of Advanced Mathematics
%%%%%%%%%%%%%%%%%%%%%%%%%%%%%%%%%%%%%%%%%%%%%%%%%%%%%%%%%%%%

\section{Linear Programming}
\label{sec:linear-programming}

\prerequisite{
Optimization fundamentals, systems of linear equations and inequalities,
matrix notation, vectors, linear algebra, convex sets, feasible regions,
and basic mathematical modeling.
}

\learningobjectives{
By the end of this section, the reader should be able to:
\begin{itemize}
    \item define a linear programming problem;
    \item distinguish objective coefficients, decision variables,
          constraint coefficients, and right-hand-side values;
    \item write linear programs in algebraic and matrix form;
    \item identify feasible, infeasible, bounded, and unbounded problems;
    \item understand geometric interpretations of linear programs;
    \item identify extreme points and vertices of feasible polyhedra;
    \item understand why linear-programming optima occur at extreme points
          when an optimum exists under appropriate conditions;
    \item convert inequalities into equality form using slack variables;
    \item distinguish slack and surplus variables;
    \item understand standard and canonical forms;
    \item identify basic and nonbasic variables;
    \item define basic solutions and basic feasible solutions;
    \item recognize degeneracy;
    \item understand the simplex method conceptually;
    \item perform elementary simplex pivots;
    \item select entering and leaving variables;
    \item use the ratio test;
    \item recognize optimality and unboundedness in simplex tableaux;
    \item understand artificial variables;
    \item understand Phase I and Phase II methods;
    \item understand the Big-\(M\) method conceptually;
    \item formulate the dual of a linear program;
    \item understand weak duality;
    \item understand strong duality;
    \item understand complementary slackness;
    \item use duality to certify optimality;
    \item interpret dual variables as shadow prices;
    \item understand reduced costs;
    \item understand sensitivity analysis;
    \item recognize alternate optimal solutions;
    \item recognize infeasibility certificates conceptually;
    \item understand Farkas' lemma conceptually;
    \item formulate transportation and blending models;
    \item recognize network-flow structure;
    \item understand interior-point methods conceptually;
    \item compare simplex and interior-point approaches;
    \item connect linear programming with convex optimization,
          operations research, economics, logistics, and computation.
\end{itemize}
}

%%%%%%%%%%%%%%%%%%%%%%%%%%%%%%%%%%%%%%%%%%%%%%%%%%%%%%%%%%%%
\subsection{What Is Linear Programming?}
\label{subsec:what-is-linear-programming}
%%%%%%%%%%%%%%%%%%%%%%%%%%%%%%%%%%%%%%%%%%%%%%%%%%%%%%%%%%%%

Linear programming, abbreviated LP, is optimization in which both the
objective function and all constraints are linear.

A standard linear program has form

\[
\boxed{
\begin{aligned}
\max_{\mathbf{x}}
\quad&
\mathbf{c}^T\mathbf{x}
\\
\text{subject to}
\quad&
A\mathbf{x}\leq\mathbf{b},
\\
&
\mathbf{x}\geq\mathbf{0}.
\end{aligned}
}
\]

Here:

\[
\mathbf{x}
=
\text{decision vector},
\]

\[
\mathbf{c}
=
\text{objective coefficients},
\]

\[
A
=
\text{constraint matrix},
\]

and

\[
\mathbf{b}
=
\text{resource or right-hand-side vector}.
\]

%%%%%%%%%%%%%%%%%%%%%%%%%%%%%%%%%%%%%%%%%%%%%%%%%%%%%%%%%%%%
\subsection{Linearity}
\label{subsec:linear-programming-linearity}
%%%%%%%%%%%%%%%%%%%%%%%%%%%%%%%%%%%%%%%%%%%%%%%%%%%%%%%%%%%%

A linear objective has form

\[
\boxed{
c_1x_1+\cdots+c_nx_n.
}
\]

A linear constraint has form

\[
\boxed{
a_1x_1+\cdots+a_nx_n
\leq
b,
}
\]

or the corresponding equality or greater-than inequality.

Products such as

\[
x_1x_2
\]

or powers such as

\[
x_1^2
\]

make the model nonlinear.

%%%%%%%%%%%%%%%%%%%%%%%%%%%%%%%%%%%%%%%%%%%%%%%%%%%%%%%%%%%%
\subsection{General Linear Program}
\label{subsec:general-linear-program}
%%%%%%%%%%%%%%%%%%%%%%%%%%%%%%%%%%%%%%%%%%%%%%%%%%%%%%%%%%%%

A general LP can be written

\[
\boxed{
\begin{aligned}
\text{optimize}
\quad&
\mathbf{c}^T\mathbf{x}
\\
\text{subject to}
\quad&
A_1\mathbf{x}\leq\mathbf{b}_1,
\\
&
A_2\mathbf{x}\geq\mathbf{b}_2,
\\
&
A_3\mathbf{x}=\mathbf{b}_3,
\end{aligned}
}
\]

with some variables possibly restricted to be nonnegative, nonpositive,
or unrestricted in sign.

%%%%%%%%%%%%%%%%%%%%%%%%%%%%%%%%%%%%%%%%%%%%%%%%%%%%%%%%%%%%
\subsection{Feasible Region}
\label{subsec:lp-feasible-region}
%%%%%%%%%%%%%%%%%%%%%%%%%%%%%%%%%%%%%%%%%%%%%%%%%%%%%%%%%%%%

The feasible region is the set of all points satisfying every constraint.

For

\[
A\mathbf{x}\leq\mathbf{b},
\]

the feasible region is

\[
\boxed{
\mathcal{F}
=
\{
\mathbf{x}:
A\mathbf{x}\leq\mathbf{b}
\}.
}
\]

Because linear inequalities define half-spaces, the feasible region of an
LP is convex.

%%%%%%%%%%%%%%%%%%%%%%%%%%%%%%%%%%%%%%%%%%%%%%%%%%%%%%%%%%%%
\subsection{Polyhedra}
\label{subsec:lp-polyhedra}
%%%%%%%%%%%%%%%%%%%%%%%%%%%%%%%%%%%%%%%%%%%%%%%%%%%%%%%%%%%%

\begin{definitionbox}{Polyhedron}
A polyhedron is a set defined by finitely many linear equalities and
inequalities.
\end{definitionbox}

Thus an LP feasible region is a polyhedron.

A bounded polyhedron is called a polytope.

%%%%%%%%%%%%%%%%%%%%%%%%%%%%%%%%%%%%%%%%%%%%%%%%%%%%%%%%%%%%
\subsection{Convexity of the Feasible Region}
\label{subsec:lp-feasible-convex}
%%%%%%%%%%%%%%%%%%%%%%%%%%%%%%%%%%%%%%%%%%%%%%%%%%%%%%%%%%%%

Suppose \(\mathbf{x}\) and \(\mathbf{y}\) are feasible:

\[
A\mathbf{x}\leq\mathbf{b},
\]

\[
A\mathbf{y}\leq\mathbf{b}.
\]

For

\[
0\leq\theta\leq1,
\]

we have

\[
A
[
\theta\mathbf{x}
+
(1-\theta)\mathbf{y}
]
=
\theta A\mathbf{x}
+
(1-\theta)A\mathbf{y}.
\]

Therefore,

\[
\boxed{
A
[
\theta\mathbf{x}
+
(1-\theta)\mathbf{y}
]
\leq
\mathbf{b}.
}
\]

Hence the feasible region is convex.

%%%%%%%%%%%%%%%%%%%%%%%%%%%%%%%%%%%%%%%%%%%%%%%%%%%%%%%%%%%%
\subsection{Geometric Interpretation}
\label{subsec:lp-geometric-interpretation}
%%%%%%%%%%%%%%%%%%%%%%%%%%%%%%%%%%%%%%%%%%%%%%%%%%%%%%%%%%%%

For two decision variables, each linear inequality defines a half-plane.

The intersection of these half-planes forms the feasible region.

The objective

\[
\boxed{
z=c_1x_1+c_2x_2
}
\]

defines parallel level lines:

\[
\boxed{
c_1x_1+c_2x_2=k.
}
\]

Optimization moves this line parallel to itself until it last touches the
feasible region.

%%%%%%%%%%%%%%%%%%%%%%%%%%%%%%%%%%%%%%%%%%%%%%%%%%%%%%%%%%%%
\subsection{Extreme Points}
\label{subsec:lp-extreme-points}
%%%%%%%%%%%%%%%%%%%%%%%%%%%%%%%%%%%%%%%%%%%%%%%%%%%%%%%%%%%%

\begin{definitionbox}{Extreme Point}
A point \(\mathbf{x}\) in a convex set \(C\) is an extreme point if it
cannot be written as

\[
\boxed{
\mathbf{x}
=
\theta\mathbf{y}
+
(1-\theta)\mathbf{z},
\qquad
0<\theta<1,
}
\]

for two distinct points

\[
\mathbf{y},\mathbf{z}\in C.
\]
\end{definitionbox}

In two dimensions, extreme points are the vertices of the feasible
polygon.

%%%%%%%%%%%%%%%%%%%%%%%%%%%%%%%%%%%%%%%%%%%%%%%%%%%%%%%%%%%%
\subsection{Fundamental LP Geometry}
\label{subsec:fundamental-lp-geometry}
%%%%%%%%%%%%%%%%%%%%%%%%%%%%%%%%%%%%%%%%%%%%%%%%%%%%%%%%%%%%

For a nonempty polyhedral feasible set, if a linear objective attains a
finite optimum, then there exists an optimal extreme point.

Thus linear programming can focus on vertices rather than every point of
the feasible region.

This geometric fact motivates the simplex method.

%%%%%%%%%%%%%%%%%%%%%%%%%%%%%%%%%%%%%%%%%%%%%%%%%%%%%%%%%%%%
\subsection{Worked Example: Graphical Linear Program}
\label{subsec:worked-graphical-lp}
%%%%%%%%%%%%%%%%%%%%%%%%%%%%%%%%%%%%%%%%%%%%%%%%%%%%%%%%%%%%

\begin{workedexamplebox}{Two-Variable Linear Program}

Maximize

\[
z=3x+2y
\]

subject to

\[
x+y\leq4,
\]

\[
x\leq2,
\]

\[
y\leq3,
\]

\[
x,y\geq0.
\]

The feasible vertices are

\[
(0,0),
\]

\[
(2,0),
\]

\[
(2,2),
\]

\[
(1,3),
\]

and

\[
(0,3).
\]

Evaluate the objective:

\[
z(0,0)=0,
\]

\[
z(2,0)=6,
\]

\[
z(2,2)=10,
\]

\[
z(1,3)=9,
\]

\[
z(0,3)=6.
\]

Therefore,

\begin{answerbox}
\[
\boxed{
(x^*,y^*)=(2,2),
}
\]

with

\[
\boxed{
z^*=10.
}
\]
\end{answerbox}

\end{workedexamplebox}

%%%%%%%%%%%%%%%%%%%%%%%%%%%%%%%%%%%%%%%%%%%%%%%%%%%%%%%%%%%%
\subsection{Bounded Linear Programs}
\label{subsec:bounded-linear-programs}
%%%%%%%%%%%%%%%%%%%%%%%%%%%%%%%%%%%%%%%%%%%%%%%%%%%%%%%%%%%%

An LP is bounded in the objective sense if its objective cannot improve
without limit over the feasible set.

For a maximization problem,

\[
\boxed{
\mathbf{c}^T\mathbf{x}
\leq M
}
\]

for some finite \(M\) over all feasible points.

A feasible set itself may be unbounded while the objective remains
bounded.

%%%%%%%%%%%%%%%%%%%%%%%%%%%%%%%%%%%%%%%%%%%%%%%%%%%%%%%%%%%%
\subsection{Unbounded Linear Programs}
\label{subsec:unbounded-linear-programs}
%%%%%%%%%%%%%%%%%%%%%%%%%%%%%%%%%%%%%%%%%%%%%%%%%%%%%%%%%%%%

A maximization LP is unbounded if feasible points exist with objective
values tending to

\[
\boxed{
+\infty.
}
\]

Similarly, a minimization LP is unbounded below if objective values tend to

\[
\boxed{
-\infty.
}
\]

Unboundedness does not mean the feasible set is empty.

%%%%%%%%%%%%%%%%%%%%%%%%%%%%%%%%%%%%%%%%%%%%%%%%%%%%%%%%%%%%
\subsection{Infeasible Linear Programs}
\label{subsec:infeasible-linear-programs}
%%%%%%%%%%%%%%%%%%%%%%%%%%%%%%%%%%%%%%%%%%%%%%%%%%%%%%%%%%%%

An LP is infeasible if no point satisfies all constraints simultaneously.

For example,

\[
x\leq1
\]

and

\[
x\geq3
\]

cannot both hold.

Thus,

\[
\boxed{
\mathcal{F}=\varnothing.
}
\]

%%%%%%%%%%%%%%%%%%%%%%%%%%%%%%%%%%%%%%%%%%%%%%%%%%%%%%%%%%%%
\subsection{Multiple Optimal Solutions}
\label{subsec:multiple-lp-optima}
%%%%%%%%%%%%%%%%%%%%%%%%%%%%%%%%%%%%%%%%%%%%%%%%%%%%%%%%%%%%

A linear program may have more than one optimal solution.

If two distinct feasible points are optimal, then every convex
combination between them is also optimal because the objective is linear.

Thus an entire edge or face may be optimal.

%%%%%%%%%%%%%%%%%%%%%%%%%%%%%%%%%%%%%%%%%%%%%%%%%%%%%%%%%%%%
\subsection{Standard Form}
\label{subsec:lp-standard-form}
%%%%%%%%%%%%%%%%%%%%%%%%%%%%%%%%%%%%%%%%%%%%%%%%%%%%%%%%%%%%

One common standard form is

\[
\boxed{
\begin{aligned}
\min_{\mathbf{x}}
\quad&
\mathbf{c}^T\mathbf{x}
\\
\text{subject to}
\quad&
A\mathbf{x}
=
\mathbf{b},
\\
&
\mathbf{x}\geq0.
\end{aligned}
}
\]

Other texts use a maximization convention.

The exact convention should always be stated.

%%%%%%%%%%%%%%%%%%%%%%%%%%%%%%%%%%%%%%%%%%%%%%%%%%%%%%%%%%%%
\subsection{Slack Variables}
\label{subsec:slack-variables}
%%%%%%%%%%%%%%%%%%%%%%%%%%%%%%%%%%%%%%%%%%%%%%%%%%%%%%%%%%%%

A constraint

\[
a_1x_1+\cdots+a_nx_n\leq b
\]

can be converted to equality by adding a nonnegative slack variable:

\[
\boxed{
a_1x_1+\cdots+a_nx_n+s=b,
}
\]

with

\[
\boxed{
s\geq0.
}
\]

The slack measures unused resource.

%%%%%%%%%%%%%%%%%%%%%%%%%%%%%%%%%%%%%%%%%%%%%%%%%%%%%%%%%%%%
\subsection{Surplus Variables}
\label{subsec:surplus-variables}
%%%%%%%%%%%%%%%%%%%%%%%%%%%%%%%%%%%%%%%%%%%%%%%%%%%%%%%%%%%%

For a constraint

\[
a_1x_1+\cdots+a_nx_n\geq b,
\]

subtract a nonnegative surplus variable:

\[
\boxed{
a_1x_1+\cdots+a_nx_n-s=b,
}
\]

where

\[
\boxed{
s\geq0.
}
\]

A surplus variable measures how much the left side exceeds the required
minimum.

%%%%%%%%%%%%%%%%%%%%%%%%%%%%%%%%%%%%%%%%%%%%%%%%%%%%%%%%%%%%
\subsection{Unrestricted Variables}
\label{subsec:unrestricted-lp-variables}
%%%%%%%%%%%%%%%%%%%%%%%%%%%%%%%%%%%%%%%%%%%%%%%%%%%%%%%%%%%%

If a variable \(x_j\) is unrestricted in sign, write

\[
\boxed{
x_j
=
x_j^+
-
x_j^-,
}
\]

where

\[
\boxed{
x_j^+\geq0,
\qquad
x_j^-\geq0.
}
\]

Thus any unrestricted variable can be represented using two nonnegative
variables.

%%%%%%%%%%%%%%%%%%%%%%%%%%%%%%%%%%%%%%%%%%%%%%%%%%%%%%%%%%%%
\subsection{Basic Solutions}
\label{subsec:lp-basic-solutions}
%%%%%%%%%%%%%%%%%%%%%%%%%%%%%%%%%%%%%%%%%%%%%%%%%%%%%%%%%%%%

Consider

\[
A\mathbf{x}=\mathbf{b},
\]

with

\[
A\in\mathbb{R}^{m\times n},
\qquad
m<n,
\]

and assume \(\operatorname{rank}(A)=m\).

Choose \(m\) linearly independent columns of \(A\).

The corresponding variables are called basic variables.

Set the remaining variables to zero and solve for the basic variables.

The resulting vector is a basic solution.

%%%%%%%%%%%%%%%%%%%%%%%%%%%%%%%%%%%%%%%%%%%%%%%%%%%%%%%%%%%%
\subsection{Basic Feasible Solution}
\label{subsec:basic-feasible-solution}
%%%%%%%%%%%%%%%%%%%%%%%%%%%%%%%%%%%%%%%%%%%%%%%%%%%%%%%%%%%%

A basic solution satisfying

\[
\boxed{
\mathbf{x}\geq0
}
\]

is a basic feasible solution, abbreviated BFS.

Under standard assumptions, basic feasible solutions correspond to
extreme points of the LP feasible polyhedron.

%%%%%%%%%%%%%%%%%%%%%%%%%%%%%%%%%%%%%%%%%%%%%%%%%%%%%%%%%%%%
\subsection{Basis Matrix}
\label{subsec:lp-basis-matrix}
%%%%%%%%%%%%%%%%%%%%%%%%%%%%%%%%%%%%%%%%%%%%%%%%%%%%%%%%%%%%

Let \(B\) be the \(m\times m\) matrix formed from basic columns.

Let \(N\) contain the nonbasic columns.

Partition

\[
A
=
[B\;N].
\]

Then

\[
B\mathbf{x}_B
+
N\mathbf{x}_N
=
\mathbf{b}.
\]

Setting

\[
\mathbf{x}_N=\mathbf{0}
\]

gives

\[
\boxed{
\mathbf{x}_B
=
B^{-1}\mathbf{b}.
}
\]

%%%%%%%%%%%%%%%%%%%%%%%%%%%%%%%%%%%%%%%%%%%%%%%%%%%%%%%%%%%%
\subsection{Degeneracy}
\label{subsec:lp-degeneracy}
%%%%%%%%%%%%%%%%%%%%%%%%%%%%%%%%%%%%%%%%%%%%%%%%%%%%%%%%%%%%

A basic feasible solution is degenerate if at least one basic variable is
zero.

Thus,

\[
\boxed{
x_{B_i}=0
}
\]

for some basic variable.

Degeneracy can cause simplex pivots that do not improve the objective.

%%%%%%%%%%%%%%%%%%%%%%%%%%%%%%%%%%%%%%%%%%%%%%%%%%%%%%%%%%%%
\subsection{The Simplex Method}
\label{subsec:simplex-method}
%%%%%%%%%%%%%%%%%%%%%%%%%%%%%%%%%%%%%%%%%%%%%%%%%%%%%%%%%%%%

The simplex method moves from one basic feasible solution to an adjacent
basic feasible solution while improving or maintaining the objective.

Conceptually:

\[
\boxed{
\text{vertex}
\rightarrow
\text{adjacent vertex}
\rightarrow
\text{better objective}.
}
\]

The process continues until an optimality condition is satisfied or
unboundedness is detected.

%%%%%%%%%%%%%%%%%%%%%%%%%%%%%%%%%%%%%%%%%%%%%%%%%%%%%%%%%%%%
\subsection{Simplex Dictionary}
\label{subsec:simplex-dictionary}
%%%%%%%%%%%%%%%%%%%%%%%%%%%%%%%%%%%%%%%%%%%%%%%%%%%%%%%%%%%%

Suppose

\[
A\mathbf{x}=\mathbf{b}
\]

is partitioned into basic and nonbasic variables.

Then

\[
B\mathbf{x}_B
+
N\mathbf{x}_N
=
\mathbf{b}.
\]

Hence,

\[
\boxed{
\mathbf{x}_B
=
B^{-1}\mathbf{b}
-
B^{-1}N\mathbf{x}_N.
}
\]

The objective becomes

\[
\boxed{
z
=
\mathbf{c}_B^TB^{-1}\mathbf{b}
+
\overline{\mathbf{c}}_N^T
\mathbf{x}_N,
}
\]

where the reduced costs depend on the chosen sign convention.

%%%%%%%%%%%%%%%%%%%%%%%%%%%%%%%%%%%%%%%%%%%%%%%%%%%%%%%%%%%%
\subsection{Reduced Costs}
\label{subsec:reduced-costs}
%%%%%%%%%%%%%%%%%%%%%%%%%%%%%%%%%%%%%%%%%%%%%%%%%%%%%%%%%%%%

For a minimization LP in standard form, a common reduced-cost formula is

\[
\boxed{
\overline{\mathbf{c}}_N
=
\mathbf{c}_N
-
N^TB^{-T}\mathbf{c}_B.
}
\]

Equivalently, if

\[
\boxed{
\mathbf{y}
=
B^{-T}\mathbf{c}_B,
}
\]

then

\[
\boxed{
\overline{\mathbf{c}}_N
=
\mathbf{c}_N
-
N^T\mathbf{y}.
}
\]

Reduced costs measure the first-order effect of increasing nonbasic
variables from zero.

%%%%%%%%%%%%%%%%%%%%%%%%%%%%%%%%%%%%%%%%%%%%%%%%%%%%%%%%%%%%
\subsection{Entering Variable}
\label{subsec:simplex-entering-variable}
%%%%%%%%%%%%%%%%%%%%%%%%%%%%%%%%%%%%%%%%%%%%%%%%%%%%%%%%%%%%

In a simplex iteration, a nonbasic variable is selected to enter the
basis.

The choice depends on the reduced-cost convention.

For a maximization tableau written in the common form, one chooses a
variable whose increase improves the objective.

Different pivot rules may choose among several candidates.

%%%%%%%%%%%%%%%%%%%%%%%%%%%%%%%%%%%%%%%%%%%%%%%%%%%%%%%%%%%%
\subsection{Leaving Variable}
\label{subsec:simplex-leaving-variable}
%%%%%%%%%%%%%%%%%%%%%%%%%%%%%%%%%%%%%%%%%%%%%%%%%%%%%%%%%%%%

As the entering variable increases, some basic variables decrease.

The leaving variable is the first basic variable to reach zero.

This maintains feasibility.

%%%%%%%%%%%%%%%%%%%%%%%%%%%%%%%%%%%%%%%%%%%%%%%%%%%%%%%%%%%%
\subsection{Ratio Test}
\label{subsec:simplex-ratio-test}
%%%%%%%%%%%%%%%%%%%%%%%%%%%%%%%%%%%%%%%%%%%%%%%%%%%%%%%%%%%%

Suppose increasing \(x_j\) changes the basic variables according to

\[
\mathbf{x}_B
=
\mathbf{b}'
-
\mathbf{d}x_j.
\]

For components with

\[
d_i>0,
\]

feasibility requires

\[
x_j
\leq
\frac{
b'_i
}{
d_i
}.
\]

Therefore the maximum feasible increase is determined by

\[
\boxed{
\min_{d_i>0}
\frac{
b'_i
}{
d_i
}.
}
\]

This is the ratio test.

%%%%%%%%%%%%%%%%%%%%%%%%%%%%%%%%%%%%%%%%%%%%%%%%%%%%%%%%%%%%
\subsection{Pivot Operation}
\label{subsec:simplex-pivot}
%%%%%%%%%%%%%%%%%%%%%%%%%%%%%%%%%%%%%%%%%%%%%%%%%%%%%%%%%%%%

A simplex pivot exchanges one basic variable and one nonbasic variable.

Algebraically, the equations are rewritten so the entering variable
becomes basic and the leaving variable becomes nonbasic.

In tableau form, this is implemented using row operations.

%%%%%%%%%%%%%%%%%%%%%%%%%%%%%%%%%%%%%%%%%%%%%%%%%%%%%%%%%%%%
\subsection{Worked Example: Slack Form}
\label{subsec:worked-lp-slack-form}
%%%%%%%%%%%%%%%%%%%%%%%%%%%%%%%%%%%%%%%%%%%%%%%%%%%%%%%%%%%%

\begin{workedexamplebox}{Introducing Slack Variables}

Consider

\[
\max
\quad
3x_1+2x_2
\]

subject to

\[
x_1+x_2\leq4,
\]

\[
x_1\leq2,
\]

\[
x_2\leq3,
\]

\[
x_1,x_2\geq0.
\]

Introduce slack variables:

\[
x_1+x_2+s_1=4,
\]

\[
x_1+s_2=2,
\]

\[
x_2+s_3=3.
\]

Thus,

\begin{answerbox}
\[
\boxed{
\begin{aligned}
x_1+x_2+s_1&=4,\\
x_1+s_2&=2,\\
x_2+s_3&=3,
\end{aligned}
}
\]

with

\[
\boxed{
x_1,x_2,s_1,s_2,s_3\geq0.
}
\]
\end{answerbox}

Setting

\[
x_1=x_2=0
\]

gives the initial basic feasible solution

\[
\boxed{
s_1=4,
\qquad
s_2=2,
\qquad
s_3=3.
}
\]

\end{workedexamplebox}

%%%%%%%%%%%%%%%%%%%%%%%%%%%%%%%%%%%%%%%%%%%%%%%%%%%%%%%%%%%%
\subsection{Simplex Optimality}
\label{subsec:simplex-optimality}
%%%%%%%%%%%%%%%%%%%%%%%%%%%%%%%%%%%%%%%%%%%%%%%%%%%%%%%%%%%%

For a minimization LP in standard form, if every nonbasic reduced cost
satisfies

\[
\boxed{
\overline c_j\geq0,
}
\]

then no feasible infinitesimal increase of a nonbasic variable can reduce
the objective.

Under the usual simplex assumptions, the current BFS is optimal.

For maximization formulations, the sign convention changes accordingly.

%%%%%%%%%%%%%%%%%%%%%%%%%%%%%%%%%%%%%%%%%%%%%%%%%%%%%%%%%%%%
\subsection{Simplex Detection of Unboundedness}
\label{subsec:simplex-unboundedness}
%%%%%%%%%%%%%%%%%%%%%%%%%%%%%%%%%%%%%%%%%%%%%%%%%%%%%%%%%%%%

Suppose a nonbasic variable can improve the objective.

If increasing it causes no basic variable to decrease to zero, then there
is no finite ratio-test bound.

The objective can improve indefinitely.

Thus the LP is unbounded in that direction.

%%%%%%%%%%%%%%%%%%%%%%%%%%%%%%%%%%%%%%%%%%%%%%%%%%%%%%%%%%%%
\subsection{Cycling}
\label{subsec:simplex-cycling}
%%%%%%%%%%%%%%%%%%%%%%%%%%%%%%%%%%%%%%%%%%%%%%%%%%%%%%%%%%%%

Degeneracy can theoretically cause the simplex method to revisit the same
basis repeatedly.

This is called cycling.

Anti-cycling pivot rules prevent this behavior.

%%%%%%%%%%%%%%%%%%%%%%%%%%%%%%%%%%%%%%%%%%%%%%%%%%%%%%%%%%%%
\subsection{Bland's Rule}
\label{subsec:blands-rule}
%%%%%%%%%%%%%%%%%%%%%%%%%%%%%%%%%%%%%%%%%%%%%%%%%%%%%%%%%%%%

Bland's rule selects entering and leaving variables according to a fixed
index ordering.

It is designed to prevent cycling.

Although it may not be computationally fastest, it provides an important
theoretical guarantee of termination.

%%%%%%%%%%%%%%%%%%%%%%%%%%%%%%%%%%%%%%%%%%%%%%%%%%%%%%%%%%%%
\subsection{Artificial Variables}
\label{subsec:artificial-variables}
%%%%%%%%%%%%%%%%%%%%%%%%%%%%%%%%%%%%%%%%%%%%%%%%%%%%%%%%%%%%

Some constraints do not produce an obvious initial basic feasible
solution.

For example,

\[
x_1+x_2=5
\]

has no slack variable that automatically provides an identity column.

An artificial variable may be introduced temporarily:

\[
\boxed{
x_1+x_2+a_1=5.
}
\]

Artificial variables are computational devices and must ultimately be
driven to zero.

%%%%%%%%%%%%%%%%%%%%%%%%%%%%%%%%%%%%%%%%%%%%%%%%%%%%%%%%%%%%
\subsection{Phase I Method}
\label{subsec:phase-one-simplex}
%%%%%%%%%%%%%%%%%%%%%%%%%%%%%%%%%%%%%%%%%%%%%%%%%%%%%%%%%%%%

Phase I searches for a feasible basis.

A common auxiliary problem minimizes the sum of artificial variables:

\[
\boxed{
\min
\sum_i a_i.
}
\]

If the optimal Phase I value is positive, the original LP is infeasible.

If the optimal value is zero, a feasible basis may be obtained.

%%%%%%%%%%%%%%%%%%%%%%%%%%%%%%%%%%%%%%%%%%%%%%%%%%%%%%%%%%%%
\subsection{Phase II Method}
\label{subsec:phase-two-simplex}
%%%%%%%%%%%%%%%%%%%%%%%%%%%%%%%%%%%%%%%%%%%%%%%%%%%%%%%%%%%%

Once a feasible basis is found, artificial variables are removed and the
original objective is restored.

The simplex method then continues on the original LP.

Thus:

\[
\boxed{
\text{Phase I}
=
\text{find feasibility},
}
\]

\[
\boxed{
\text{Phase II}
=
\text{optimize}.
}
\]

%%%%%%%%%%%%%%%%%%%%%%%%%%%%%%%%%%%%%%%%%%%%%%%%%%%%%%%%%%%%
\subsection{Big-\(M\) Method Preview}
\label{subsec:big-m-method}
%%%%%%%%%%%%%%%%%%%%%%%%%%%%%%%%%%%%%%%%%%%%%%%%%%%%%%%%%%%%

An alternative places a very large penalty on artificial variables.

For a minimization problem, one may use an objective containing

\[
\boxed{
M\sum_i a_i,
}
\]

where \(M\) is symbolically very large and positive.

The method encourages artificial variables to leave the basis.

Two-phase methods are often cleaner theoretically and numerically.

%%%%%%%%%%%%%%%%%%%%%%%%%%%%%%%%%%%%%%%%%%%%%%%%%%%%%%%%%%%%
\subsection{Duality}
\label{subsec:linear-programming-duality}
%%%%%%%%%%%%%%%%%%%%%%%%%%%%%%%%%%%%%%%%%%%%%%%%%%%%%%%%%%%%

Every linear program has an associated dual problem.

For the primal maximization problem

\[
\boxed{
\begin{aligned}
\max_{\mathbf{x}}
\quad&
\mathbf{c}^T\mathbf{x}
\\
\text{subject to}
\quad&
A\mathbf{x}\leq\mathbf{b},
\\
&
\mathbf{x}\geq0,
\end{aligned}
}
\]

the corresponding dual is

\[
\boxed{
\begin{aligned}
\min_{\mathbf{y}}
\quad&
\mathbf{b}^T\mathbf{y}
\\
\text{subject to}
\quad&
A^T\mathbf{y}\geq\mathbf{c},
\\
&
\mathbf{y}\geq0.
\end{aligned}
}
\]

%%%%%%%%%%%%%%%%%%%%%%%%%%%%%%%%%%%%%%%%%%%%%%%%%%%%%%%%%%%%
\subsection{Primal and Dual Variables}
\label{subsec:primal-dual-variables}
%%%%%%%%%%%%%%%%%%%%%%%%%%%%%%%%%%%%%%%%%%%%%%%%%%%%%%%%%%%%

Each primal constraint corresponds to a dual variable.

Each primal variable corresponds to a dual constraint.

Thus the structure is reversed:

\[
\boxed{
\text{primal constraints}
\longleftrightarrow
\text{dual variables},
}
\]

and

\[
\boxed{
\text{primal variables}
\longleftrightarrow
\text{dual constraints}.
}
\]

%%%%%%%%%%%%%%%%%%%%%%%%%%%%%%%%%%%%%%%%%%%%%%%%%%%%%%%%%%%%
\subsection{Weak Duality}
\label{subsec:weak-duality}
%%%%%%%%%%%%%%%%%%%%%%%%%%%%%%%%%%%%%%%%%%%%%%%%%%%%%%%%%%%%

Suppose \(\mathbf{x}\) is primal feasible and \(\mathbf{y}\) is dual
feasible.

Then

\[
A\mathbf{x}\leq\mathbf{b},
\]

and

\[
A^T\mathbf{y}\geq\mathbf{c}.
\]

Because

\[
\mathbf{x}\geq0,
\qquad
\mathbf{y}\geq0,
\]

we obtain

\[
\mathbf{c}^T\mathbf{x}
\leq
\mathbf{y}^TA\mathbf{x}
\leq
\mathbf{y}^T\mathbf{b}.
\]

Therefore,

\[
\boxed{
\mathbf{c}^T\mathbf{x}
\leq
\mathbf{b}^T\mathbf{y}.
}
\]

This is weak duality.

%%%%%%%%%%%%%%%%%%%%%%%%%%%%%%%%%%%%%%%%%%%%%%%%%%%%%%%%%%%%
\subsection{Interpretation of Weak Duality}
\label{subsec:weak-duality-interpretation}
%%%%%%%%%%%%%%%%%%%%%%%%%%%%%%%%%%%%%%%%%%%%%%%%%%%%%%%%%%%%

For a primal maximization problem:

\[
\boxed{
\text{every dual feasible value is an upper bound on primal profit}.
}
\]

For a dual minimization problem:

\[
\boxed{
\text{every primal feasible value is a lower bound on dual cost}.
}
\]

If the values become equal, both solutions must be optimal.

%%%%%%%%%%%%%%%%%%%%%%%%%%%%%%%%%%%%%%%%%%%%%%%%%%%%%%%%%%%%
\subsection{Strong Duality}
\label{subsec:strong-duality}
%%%%%%%%%%%%%%%%%%%%%%%%%%%%%%%%%%%%%%%%%%%%%%%%%%%%%%%%%%%%

If the primal and dual LPs have finite optimal solutions, then

\[
\boxed{
\mathbf{c}^T\mathbf{x}^*
=
\mathbf{b}^T\mathbf{y}^*.
}
\]

This is the strong duality theorem.

Thus there is no duality gap for ordinary finite-dimensional linear
programming under the usual feasibility and boundedness assumptions.

%%%%%%%%%%%%%%%%%%%%%%%%%%%%%%%%%%%%%%%%%%%%%%%%%%%%%%%%%%%%
\subsection{Duality Gap}
\label{subsec:lp-duality-gap}
%%%%%%%%%%%%%%%%%%%%%%%%%%%%%%%%%%%%%%%%%%%%%%%%%%%%%%%%%%%%

For primal feasible \(\mathbf{x}\) and dual feasible \(\mathbf{y}\), define

\[
\boxed{
\mathbf{b}^T\mathbf{y}
-
\mathbf{c}^T\mathbf{x}.
}
\]

Weak duality implies this gap is nonnegative.

At optimality,

\[
\boxed{
\mathbf{b}^T\mathbf{y}^*
-
\mathbf{c}^T\mathbf{x}^*
=
0.
}
\]

%%%%%%%%%%%%%%%%%%%%%%%%%%%%%%%%%%%%%%%%%%%%%%%%%%%%%%%%%%%%
\subsection{Worked Example: Constructing the Dual}
\label{subsec:worked-lp-dual}
%%%%%%%%%%%%%%%%%%%%%%%%%%%%%%%%%%%%%%%%%%%%%%%%%%%%%%%%%%%%

\begin{workedexamplebox}{Primal to Dual}

Consider

\[
\begin{aligned}
\max
\quad&
3x_1+2x_2
\\
\text{subject to}
\quad&
x_1+x_2\leq4,
\\
&
x_1\leq2,
\\
&
x_2\leq3,
\\
&
x_1,x_2\geq0.
\end{aligned}
\]

Introduce dual variables

\[
y_1,y_2,y_3\geq0.
\]

The dual objective is

\[
4y_1+2y_2+3y_3.
\]

The dual constraints are obtained column by column:

\[
y_1+y_2\geq3,
\]

\[
y_1+y_3\geq2.
\]

Therefore,

\begin{answerbox}
\[
\boxed{
\begin{aligned}
\min
\quad&
4y_1+2y_2+3y_3
\\
\text{subject to}
\quad&
y_1+y_2\geq3,
\\
&
y_1+y_3\geq2,
\\
&
y_1,y_2,y_3\geq0.
\end{aligned}
}
\]
\end{answerbox}

\end{workedexamplebox}

%%%%%%%%%%%%%%%%%%%%%%%%%%%%%%%%%%%%%%%%%%%%%%%%%%%%%%%%%%%%
\subsection{Complementary Slackness}
\label{subsec:lp-complementary-slackness}
%%%%%%%%%%%%%%%%%%%%%%%%%%%%%%%%%%%%%%%%%%%%%%%%%%%%%%%%%%%%

For the standard primal-dual pair,

\[
\max
\{
\mathbf{c}^T\mathbf{x}:
A\mathbf{x}\leq\mathbf{b},
\mathbf{x}\geq0
\},
\]

and

\[
\min
\{
\mathbf{b}^T\mathbf{y}:
A^T\mathbf{y}\geq\mathbf{c},
\mathbf{y}\geq0
\},
\]

optimal solutions satisfy

\[
\boxed{
y_i
[
b_i-(A\mathbf{x})_i
]
=
0,
}
\]

and

\[
\boxed{
x_j
[
(A^T\mathbf{y})_j-c_j
]
=
0.
}
\]

%%%%%%%%%%%%%%%%%%%%%%%%%%%%%%%%%%%%%%%%%%%%%%%%%%%%%%%%%%%%
\subsection{Complementary Slackness Interpretation}
\label{subsec:complementary-slackness-interpretation}
%%%%%%%%%%%%%%%%%%%%%%%%%%%%%%%%%%%%%%%%%%%%%%%%%%%%%%%%%%%%

If primal constraint \(i\) has slack,

\[
b_i-(A\mathbf{x})_i>0,
\]

then

\[
\boxed{
y_i=0.
}
\]

If primal variable \(x_j\) is positive, then the corresponding dual
constraint must be tight:

\[
\boxed{
(A^T\mathbf{y})_j=c_j.
}
\]

%%%%%%%%%%%%%%%%%%%%%%%%%%%%%%%%%%%%%%%%%%%%%%%%%%%%%%%%%%%%
\subsection{Worked Example: Complementary Slackness}
\label{subsec:worked-complementary-slackness}
%%%%%%%%%%%%%%%%%%%%%%%%%%%%%%%%%%%%%%%%%%%%%%%%%%%%%%%%%%%%

\begin{workedexamplebox}{Recovering Dual Information}

For the earlier primal optimum

\[
(x_1,x_2)=(2,2),
\]

the constraints satisfy

\[
x_1+x_2=4,
\]

\[
x_1=2,
\]

while

\[
x_2=2<3.
\]

Thus the third primal constraint has positive slack.

Complementary slackness gives

\[
\boxed{
y_3=0.
}
\]

Since

\[
x_1>0,
\qquad
x_2>0,
\]

the dual constraints must both be tight:

\[
y_1+y_2=3,
\]

\[
y_1+y_3=2.
\]

Since \(y_3=0\),

\[
y_1=2,
\]

and therefore

\[
y_2=1.
\]

The dual objective is

\[
4(2)+2(1)+3(0)=10.
\]

Hence,

\begin{answerbox}
\[
\boxed{
\mathbf{y}^*
=
(2,1,0),
}
\]

and

\[
\boxed{
\mathbf{b}^T\mathbf{y}^*
=
10
=
\mathbf{c}^T\mathbf{x}^*.
}
\]
\end{answerbox}

\end{workedexamplebox}

%%%%%%%%%%%%%%%%%%%%%%%%%%%%%%%%%%%%%%%%%%%%%%%%%%%%%%%%%%%%
\subsection{Dual Variables as Shadow Prices}
\label{subsec:dual-shadow-prices}
%%%%%%%%%%%%%%%%%%%%%%%%%%%%%%%%%%%%%%%%%%%%%%%%%%%%%%%%%%%%

A dual variable often measures the marginal value of increasing a resource
limit.

If the \(i\)-th resource constraint is

\[
a_i^T x\leq b_i,
\]

then under a stable optimal basis,

\[
\boxed{
y_i^*
}
\]

represents the approximate change in the optimal objective per unit
increase in \(b_i\).

Thus dual variables are often called shadow prices.

%%%%%%%%%%%%%%%%%%%%%%%%%%%%%%%%%%%%%%%%%%%%%%%%%%%%%%%%%%%%
\subsection{Sensitivity to Right-Hand-Side Changes}
\label{subsec:lp-rhs-sensitivity}
%%%%%%%%%%%%%%%%%%%%%%%%%%%%%%%%%%%%%%%%%%%%%%%%%%%%%%%%%%%%

Suppose

\[
\mathbf{b}
\]

changes to

\[
\mathbf{b}+\Delta\mathbf{b}.
\]

If the same basis remains optimal, then the basic solution becomes

\[
\boxed{
\mathbf{x}_B
=
B^{-1}
(
\mathbf{b}+\Delta\mathbf{b}
).
}
\]

Feasibility requires

\[
\boxed{
B^{-1}
(
\mathbf{b}+\Delta\mathbf{b}
)
\geq0.
}
\]

This determines allowable ranges for resource perturbations.

%%%%%%%%%%%%%%%%%%%%%%%%%%%%%%%%%%%%%%%%%%%%%%%%%%%%%%%%%%%%
\subsection{Sensitivity of the Optimal Value}
\label{subsec:lp-optimal-value-sensitivity}
%%%%%%%%%%%%%%%%%%%%%%%%%%%%%%%%%%%%%%%%%%%%%%%%%%%%%%%%%%%%

When the optimal basis remains unchanged, a small right-hand-side change
has objective effect

\[
\boxed{
\Delta z^*
\approx
(\mathbf{y}^*)^T
\Delta\mathbf{b}.
}
\]

Within the valid sensitivity range, this relation can be exact for an LP.

%%%%%%%%%%%%%%%%%%%%%%%%%%%%%%%%%%%%%%%%%%%%%%%%%%%%%%%%%%%%
\subsection{Objective-Coefficient Sensitivity}
\label{subsec:lp-objective-sensitivity}
%%%%%%%%%%%%%%%%%%%%%%%%%%%%%%%%%%%%%%%%%%%%%%%%%%%%%%%%%%%%

Changes in

\[
\mathbf{c}
\]

alter the slope of the objective but not the feasible region.

The current basis remains optimal while the associated reduced costs
retain the required optimality signs.

Thus sensitivity analysis can determine ranges of objective coefficients
that preserve the optimal basis.

%%%%%%%%%%%%%%%%%%%%%%%%%%%%%%%%%%%%%%%%%%%%%%%%%%%%%%%%%%%%
\subsection{Reduced-Cost Interpretation}
\label{subsec:lp-reduced-cost-interpretation}
%%%%%%%%%%%%%%%%%%%%%%%%%%%%%%%%%%%%%%%%%%%%%%%%%%%%%%%%%%%%

For a minimization LP, the reduced cost

\[
\boxed{
\overline c_j
=
c_j-a_j^Ty
}
\]

measures the net objective cost of increasing nonbasic variable \(x_j\).

At an optimal basis,

\[
\boxed{
\overline c_j\geq0.
}
\]

For a maximization formulation, the corresponding sign convention is
reversed.

%%%%%%%%%%%%%%%%%%%%%%%%%%%%%%%%%%%%%%%%%%%%%%%%%%%%%%%%%%%%
\subsection{Alternate Optimal Solutions in Simplex}
\label{subsec:simplex-alternate-optimum}
%%%%%%%%%%%%%%%%%%%%%%%%%%%%%%%%%%%%%%%%%%%%%%%%%%%%%%%%%%%%

At an optimal BFS, if a nonbasic variable has zero reduced cost, then
moving it into the basis may produce another optimal solution.

Thus

\[
\boxed{
\overline c_j=0
}
\]

for a nonbasic variable can signal alternate optima.

%%%%%%%%%%%%%%%%%%%%%%%%%%%%%%%%%%%%%%%%%%%%%%%%%%%%%%%%%%%%
\subsection{Dual of the Dual}
\label{subsec:dual-of-dual}
%%%%%%%%%%%%%%%%%%%%%%%%%%%%%%%%%%%%%%%%%%%%%%%%%%%%%%%%%%%%

Under consistent standard-form conventions, taking the dual of the dual
returns the primal problem.

Thus,

\[
\boxed{
(P^*)^*=P.
}
\]

This reflects the structural symmetry of LP duality.

%%%%%%%%%%%%%%%%%%%%%%%%%%%%%%%%%%%%%%%%%%%%%%%%%%%%%%%%%%%%
\subsection{Constraint and Variable Sign Rules}
\label{subsec:lp-duality-sign-rules}
%%%%%%%%%%%%%%%%%%%%%%%%%%%%%%%%%%%%%%%%%%%%%%%%%%%%%%%%%%%%

Duality rules depend on the chosen primal orientation.

For a maximization primal:

\[
\boxed{
\text{primal } \leq
\text{ constraint}
\leftrightarrow
\text{nonnegative dual variable},
}
\]

\[
\boxed{
\text{primal equality}
\leftrightarrow
\text{unrestricted dual variable}.
}
\]

Similarly, variable sign restrictions affect the corresponding dual
constraint direction.

A dual should therefore be derived systematically rather than memorized
from one special case only.

%%%%%%%%%%%%%%%%%%%%%%%%%%%%%%%%%%%%%%%%%%%%%%%%%%%%%%%%%%%%
\subsection{Farkas' Lemma Preview}
\label{subsec:farkas-lemma}
%%%%%%%%%%%%%%%%%%%%%%%%%%%%%%%%%%%%%%%%%%%%%%%%%%%%%%%%%%%%

Farkas' lemma provides an alternative theorem for linear inequalities.

One standard version states that exactly one of the following can hold:

\[
\boxed{
A\mathbf{x}
=
\mathbf{b},
\qquad
\mathbf{x}\geq0,
}
\]

or there exists \(\mathbf{y}\) such that

\[
\boxed{
A^T\mathbf{y}\geq0,
}
\]

and

\[
\boxed{
\mathbf{b}^T\mathbf{y}<0.
}
\]

The vector \(\mathbf{y}\) acts as a certificate of infeasibility.

%%%%%%%%%%%%%%%%%%%%%%%%%%%%%%%%%%%%%%%%%%%%%%%%%%%%%%%%%%%%
\subsection{Certificates of Optimality}
\label{subsec:lp-certificates-optimality}
%%%%%%%%%%%%%%%%%%%%%%%%%%%%%%%%%%%%%%%%%%%%%%%%%%%%%%%%%%%%

Suppose one finds:

\[
\boxed{
\text{primal feasible }\mathbf{x},
}
\]

and

\[
\boxed{
\text{dual feasible }\mathbf{y},
}
\]

with equal objective values:

\[
\boxed{
\mathbf{c}^T\mathbf{x}
=
\mathbf{b}^T\mathbf{y}.
}
\]

Weak duality then proves both solutions are optimal.

This provides a powerful certificate of optimality.

%%%%%%%%%%%%%%%%%%%%%%%%%%%%%%%%%%%%%%%%%%%%%%%%%%%%%%%%%%%%
\subsection{Transportation Problems}
\label{subsec:transportation-problems}
%%%%%%%%%%%%%%%%%%%%%%%%%%%%%%%%%%%%%%%%%%%%%%%%%%%%%%%%%%%%

Suppose goods are shipped from origins \(i\) to destinations \(j\).

Let

\[
x_{ij}
=
\text{quantity shipped from }i\text{ to }j.
\]

If

\[
c_{ij}
\]

is unit transportation cost, then the objective is

\[
\boxed{
\min
\sum_i
\sum_j
c_{ij}x_{ij}.
}
\]

%%%%%%%%%%%%%%%%%%%%%%%%%%%%%%%%%%%%%%%%%%%%%%%%%%%%%%%%%%%%
\subsection{Transportation Constraints}
\label{subsec:transportation-constraints}
%%%%%%%%%%%%%%%%%%%%%%%%%%%%%%%%%%%%%%%%%%%%%%%%%%%%%%%%%%%%

If origin \(i\) has supply \(s_i\),

\[
\boxed{
\sum_j
x_{ij}
\leq
s_i.
}
\]

If destination \(j\) requires demand \(d_j\),

\[
\boxed{
\sum_i
x_{ij}
\geq
d_j.
}
\]

For a balanced transportation problem,

\[
\boxed{
\sum_i s_i
=
\sum_j d_j.
}
\]

Then supply and demand constraints are often written as equalities.

%%%%%%%%%%%%%%%%%%%%%%%%%%%%%%%%%%%%%%%%%%%%%%%%%%%%%%%%%%%%
\subsection{Blending Problems}
\label{subsec:blending-problems}
%%%%%%%%%%%%%%%%%%%%%%%%%%%%%%%%%%%%%%%%%%%%%%%%%%%%%%%%%%%%

Suppose ingredients \(x_i\) are combined to produce a mixture.

If ingredient \(i\) costs \(c_i\),

\[
\boxed{
\min
\sum_i
c_ix_i.
}
\]

Quality requirements produce linear constraints such as

\[
\boxed{
\sum_i
a_{ji}x_i
\geq
b_j.
}
\]

Applications include fuels, diets, alloys, feeds, and chemical mixtures.

%%%%%%%%%%%%%%%%%%%%%%%%%%%%%%%%%%%%%%%%%%%%%%%%%%%%%%%%%%%%
\subsection{Diet Problem}
\label{subsec:diet-problem}
%%%%%%%%%%%%%%%%%%%%%%%%%%%%%%%%%%%%%%%%%%%%%%%%%%%%%%%%%%%%

The classical diet problem minimizes food cost while meeting nutritional
requirements.

Let

\[
x_j
=
\text{quantity of food }j.
\]

Then

\[
\boxed{
\min
\sum_j
c_jx_j
}
\]

subject to

\[
\boxed{
\sum_j
a_{ij}x_j
\geq
b_i
}
\]

for each nutrient \(i\).

%%%%%%%%%%%%%%%%%%%%%%%%%%%%%%%%%%%%%%%%%%%%%%%%%%%%%%%%%%%%
\subsection{Production Planning}
\label{subsec:lp-production-planning}
%%%%%%%%%%%%%%%%%%%%%%%%%%%%%%%%%%%%%%%%%%%%%%%%%%%%%%%%%%%%

A production LP may maximize profit:

\[
\boxed{
\max
\sum_j
p_jx_j
}
\]

subject to resource limits:

\[
\boxed{
\sum_j
a_{ij}x_j
\leq
b_i.
}
\]

This is one of the standard operations-research formulations.

%%%%%%%%%%%%%%%%%%%%%%%%%%%%%%%%%%%%%%%%%%%%%%%%%%%%%%%%%%%%
\subsection{Network Flow Preview}
\label{subsec:network-flow-lp-preview}
%%%%%%%%%%%%%%%%%%%%%%%%%%%%%%%%%%%%%%%%%%%%%%%%%%%%%%%%%%%%

Many network problems are linear programs.

Let

\[
x_{ij}
\]

denote flow along arc \((i,j)\).

A flow-conservation equation at node \(i\) is

\[
\boxed{
\sum_j
x_{ji}
-
\sum_j
x_{ij}
=
b_i.
}
\]

Capacity constraints have form

\[
\boxed{
0\leq x_{ij}\leq u_{ij}.
}
\]

Network optimization is developed later in Chapter 10.

%%%%%%%%%%%%%%%%%%%%%%%%%%%%%%%%%%%%%%%%%%%%%%%%%%%%%%%%%%%%
\subsection{Integrality in Special LPs}
\label{subsec:special-lp-integrality}
%%%%%%%%%%%%%%%%%%%%%%%%%%%%%%%%%%%%%%%%%%%%%%%%%%%%%%%%%%%%

Some linear programs with integral data automatically possess integral
optimal extreme points.

This occurs for important classes involving totally unimodular matrices.

Examples include certain:

\[
\boxed{
\text{network-flow problems},
}
\]

\[
\boxed{
\text{assignment problems},
}
\]

and

\[
\boxed{
\text{transportation problems}.
}
\]

Thus some apparently discrete problems can be solved as ordinary LPs.

%%%%%%%%%%%%%%%%%%%%%%%%%%%%%%%%%%%%%%%%%%%%%%%%%%%%%%%%%%%%
\subsection{Total Unimodularity Preview}
\label{subsec:total-unimodularity-preview}
%%%%%%%%%%%%%%%%%%%%%%%%%%%%%%%%%%%%%%%%%%%%%%%%%%%%%%%%%%%%

A matrix is totally unimodular if every square subdeterminant is in

\[
\boxed{
\{-1,0,1\}.
}
\]

Under suitable conditions, if \(A\) is totally unimodular and
\(\mathbf{b}\) is integral, then vertices of

\[
\{
x:
Ax\leq b,
x\geq0
\}
\]

are integral.

This property becomes important in combinatorial optimization.

%%%%%%%%%%%%%%%%%%%%%%%%%%%%%%%%%%%%%%%%%%%%%%%%%%%%%%%%%%%%
\subsection{Interior-Point Methods}
\label{subsec:lp-interior-point-methods}
%%%%%%%%%%%%%%%%%%%%%%%%%%%%%%%%%%%%%%%%%%%%%%%%%%%%%%%%%%%%

The simplex method moves along the boundary of the feasible polyhedron.

Interior-point methods instead move through the interior.

A barrier formulation for

\[
A\mathbf{x}\leq\mathbf{b}
\]

may include logarithmic terms such as

\[
\boxed{
-\mu
\sum_i
\log
[
b_i-a_i^T\mathbf{x}
].
}
\]

As

\[
\mu\to0,
\]

solutions approach the LP optimum under suitable conditions.

%%%%%%%%%%%%%%%%%%%%%%%%%%%%%%%%%%%%%%%%%%%%%%%%%%%%%%%%%%%%
\subsection{Central Path}
\label{subsec:lp-central-path}
%%%%%%%%%%%%%%%%%%%%%%%%%%%%%%%%%%%%%%%%%%%%%%%%%%%%%%%%%%%%

Interior-point methods generate a family of solutions parameterized by a
barrier parameter \(\mu>0\).

These points trace a curve called the central path.

As

\[
\boxed{
\mu\to0,
}
\]

the central path approaches an optimal primal-dual solution.

%%%%%%%%%%%%%%%%%%%%%%%%%%%%%%%%%%%%%%%%%%%%%%%%%%%%%%%%%%%%
\subsection{Primal-Dual Interior-Point Methods}
\label{subsec:primal-dual-interior-point}
%%%%%%%%%%%%%%%%%%%%%%%%%%%%%%%%%%%%%%%%%%%%%%%%%%%%%%%%%%%%

Modern interior-point methods often solve primal and dual conditions
simultaneously.

They approximately enforce:

\[
\boxed{
\text{primal feasibility},
}
\]

\[
\boxed{
\text{dual feasibility},
}
\]

and a perturbed complementary-slackness condition such as

\[
\boxed{
x_js_j
=
\mu.
}
\]

Newton's method is then applied to the resulting nonlinear system.

%%%%%%%%%%%%%%%%%%%%%%%%%%%%%%%%%%%%%%%%%%%%%%%%%%%%%%%%%%%%
\subsection{Simplex Versus Interior-Point Methods}
\label{subsec:simplex-versus-interior-point}
%%%%%%%%%%%%%%%%%%%%%%%%%%%%%%%%%%%%%%%%%%%%%%%%%%%%%%%%%%%%

The simplex method:

\[
\boxed{
\text{moves between basic feasible solutions}.
}
\]

Interior-point methods:

\[
\boxed{
\text{move through the interior}.
}
\]

Simplex methods can be excellent for sensitivity analysis and warm starts.

Interior-point methods have strong polynomial-time theory and can be very
effective for large sparse LPs.

%%%%%%%%%%%%%%%%%%%%%%%%%%%%%%%%%%%%%%%%%%%%%%%%%%%%%%%%%%%%
\subsection{Computational Complexity}
\label{subsec:lp-computational-complexity}
%%%%%%%%%%%%%%%%%%%%%%%%%%%%%%%%%%%%%%%%%%%%%%%%%%%%%%%%%%%%

The simplex method performs extremely well in many practical problems but
has exponential worst-case examples for common pivot rules.

Polynomial-time LP algorithms include:

\[
\boxed{
\text{ellipsoid methods},
}
\]

and

\[
\boxed{
\text{interior-point methods}.
}
\]

Thus linear programming is polynomial-time solvable in a theoretical
complexity sense.

%%%%%%%%%%%%%%%%%%%%%%%%%%%%%%%%%%%%%%%%%%%%%%%%%%%%%%%%%%%%
\subsection{Presolve}
\label{subsec:lp-presolve}
%%%%%%%%%%%%%%%%%%%%%%%%%%%%%%%%%%%%%%%%%%%%%%%%%%%%%%%%%%%%

Modern solvers simplify LP models before optimization.

Presolve may:

\[
\boxed{
\text{remove redundant constraints},
}
\]

\[
\boxed{
\text{fix variables},
}
\]

\[
\boxed{
\text{tighten bounds},
}
\]

or

\[
\boxed{
\text{detect infeasibility}.
}
\]

A smaller equivalent model can be much easier to solve.

%%%%%%%%%%%%%%%%%%%%%%%%%%%%%%%%%%%%%%%%%%%%%%%%%%%%%%%%%%%%
\subsection{Scaling in Linear Programming}
\label{subsec:lp-scaling}
%%%%%%%%%%%%%%%%%%%%%%%%%%%%%%%%%%%%%%%%%%%%%%%%%%%%%%%%%%%%

Constraint coefficients with widely different magnitudes can cause
numerical difficulties.

For example, coefficients of order

\[
10^{-9}
\]

and

\[
10^9
\]

in the same model may produce poor conditioning.

Scaling rows and columns can improve numerical reliability.

%%%%%%%%%%%%%%%%%%%%%%%%%%%%%%%%%%%%%%%%%%%%%%%%%%%%%%%%%%%%
\subsection{Redundant Constraints}
\label{subsec:lp-redundant-constraints}
%%%%%%%%%%%%%%%%%%%%%%%%%%%%%%%%%%%%%%%%%%%%%%%%%%%%%%%%%%%%

A constraint is redundant if removing it does not change the feasible
region.

For example, if

\[
x\leq2
\]

is already required, then

\[
x\leq5
\]

adds no new restriction.

Redundant constraints may still influence numerical performance.

%%%%%%%%%%%%%%%%%%%%%%%%%%%%%%%%%%%%%%%%%%%%%%%%%%%%%%%%%%%%
\subsection{Binding Constraints}
\label{subsec:binding-constraints-lp}
%%%%%%%%%%%%%%%%%%%%%%%%%%%%%%%%%%%%%%%%%%%%%%%%%%%%%%%%%%%%

A constraint is binding at a feasible point if it holds with equality.

For

\[
a_i^Tx\leq b_i,
\]

the constraint is binding when

\[
\boxed{
a_i^Tx=b_i.
}
\]

Otherwise it has positive slack.

%%%%%%%%%%%%%%%%%%%%%%%%%%%%%%%%%%%%%%%%%%%%%%%%%%%%%%%%%%%%
\subsection{LP Modeling Workflow}
\label{subsec:lp-modeling-workflow}
%%%%%%%%%%%%%%%%%%%%%%%%%%%%%%%%%%%%%%%%%%%%%%%%%%%%%%%%%%%%

A practical linear-programming workflow is:

\begin{enumerate}
    \item identify the decisions;
    \item define decision variables with units;
    \item identify the quantity to minimize or maximize;
    \item write the linear objective;
    \item identify every resource, demand, logical, or policy constraint;
    \item write each constraint linearly;
    \item specify sign restrictions;
    \item check whether the feasible region is nonempty;
    \item solve the LP;
    \item verify primal feasibility;
    \item inspect binding constraints and slack;
    \item inspect dual values and reduced costs;
    \item conduct sensitivity analysis;
    \item interpret the solution in application terms.
\end{enumerate}

%%%%%%%%%%%%%%%%%%%%%%%%%%%%%%%%%%%%%%%%%%%%%%%%%%%%%%%%%%%%
\subsection{Common Errors}
\label{subsec:linear-programming-common-errors}
%%%%%%%%%%%%%%%%%%%%%%%%%%%%%%%%%%%%%%%%%%%%%%%%%%%%%%%%%%%%

\subsubsection{Using Nonlinear Expressions in an LP}

Terms such as

\[
x_1x_2,
\]

\[
x_1^2,
\]

and

\[
\frac{x_1}{x_2}
\]

are nonlinear.

A model containing them is not an ordinary linear program.

%%%%%%%%%%%%%%%%%%%%%%%%%%%%%%%%%%%%%%%%%%%%%%%%%%%%%%%%%%%%
\subsubsection{Forgetting Nonnegativity Restrictions}

If a production quantity cannot be negative, include

\[
\boxed{
x_j\geq0.
}
\]

Without that restriction, the mathematical model may allow meaningless
solutions.

%%%%%%%%%%%%%%%%%%%%%%%%%%%%%%%%%%%%%%%%%%%%%%%%%%%%%%%%%%%%
\subsubsection{Confusing Slack and Surplus Variables}

For

\[
\leq,
\]

add slack.

For

\[
\geq,
\]

subtract surplus.

The resulting auxiliary variables remain nonnegative.

%%%%%%%%%%%%%%%%%%%%%%%%%%%%%%%%%%%%%%%%%%%%%%%%%%%%%%%%%%%%
\subsubsection{Assuming Every Feasible Vertex Is Optimal}

Vertices are candidates.

The objective must still be compared or optimized.

%%%%%%%%%%%%%%%%%%%%%%%%%%%%%%%%%%%%%%%%%%%%%%%%%%%%%%%%%%%%
\subsubsection{Assuming an Unbounded Feasible Region Implies an Unbounded Objective}

A feasible region may extend infinitely while the objective still has a
finite optimum.

Unboundedness must be checked relative to the optimization direction.

%%%%%%%%%%%%%%%%%%%%%%%%%%%%%%%%%%%%%%%%%%%%%%%%%%%%%%%%%%%%
\subsubsection{Confusing Infeasible and Unbounded}

Infeasible means:

\[
\boxed{
\text{no feasible point exists}.
}
\]

Unbounded means:

\[
\boxed{
\text{feasible points exist, but the objective improves without bound}.
}
\]

These are fundamentally different outcomes.

%%%%%%%%%%%%%%%%%%%%%%%%%%%%%%%%%%%%%%%%%%%%%%%%%%%%%%%%%%%%
\subsubsection{Ignoring Degeneracy}

A basic variable may equal zero.

This can cause a pivot with no objective improvement and can complicate
simplex behavior.

%%%%%%%%%%%%%%%%%%%%%%%%%%%%%%%%%%%%%%%%%%%%%%%%%%%%%%%%%%%%
\subsubsection{Using the Ratio Test with Nonpositive Pivot Coefficients}

Only rows that actually limit the entering variable should be included in
the ratio test.

Using inappropriate signs can destroy feasibility.

%%%%%%%%%%%%%%%%%%%%%%%%%%%%%%%%%%%%%%%%%%%%%%%%%%%%%%%%%%%%
\subsubsection{Assuming Artificial Variables Are Real Decision Variables}

Artificial variables are temporary computational devices.

A feasible solution to the original LP must have them equal to zero.

%%%%%%%%%%%%%%%%%%%%%%%%%%%%%%%%%%%%%%%%%%%%%%%%%%%%%%%%%%%%
\subsubsection{Mixing Dual Sign Conventions}

Dual constraints and variable signs depend on whether the primal is
written as minimization or maximization and on inequality orientation.

Use one convention consistently.

%%%%%%%%%%%%%%%%%%%%%%%%%%%%%%%%%%%%%%%%%%%%%%%%%%%%%%%%%%%%
\subsubsection{Forgetting Weak Duality}

For a standard primal maximization and dual minimization pair,

\[
\boxed{
\text{primal feasible value}
\leq
\text{dual feasible value}.
}
\]

A violation indicates an algebraic or formulation error.

%%%%%%%%%%%%%%%%%%%%%%%%%%%%%%%%%%%%%%%%%%%%%%%%%%%%%%%%%%%%
\subsubsection{Treating Shadow Prices as Globally Valid}

A shadow price describes marginal objective sensitivity only within a
range where the relevant optimal basis or local dual structure remains
valid.

Large parameter changes can alter the active constraints.

%%%%%%%%%%%%%%%%%%%%%%%%%%%%%%%%%%%%%%%%%%%%%%%%%%%%%%%%%%%%
\subsubsection{Ignoring Numerical Tolerances}

Practical solvers use tolerances.

A reported value such as

\[
-10^{-10}
\]

for a nonnegative variable may reflect floating-point error rather than a
meaningful violation.

%%%%%%%%%%%%%%%%%%%%%%%%%%%%%%%%%%%%%%%%%%%%%%%%%%%%%%%%%%%%
\subsection{Applications}
\label{subsec:linear-programming-applications}
%%%%%%%%%%%%%%%%%%%%%%%%%%%%%%%%%%%%%%%%%%%%%%%%%%%%%%%%%%%%

\begin{applicationbox}

\textbf{Production Planning.}
Linear programs allocate labor, materials, machine time, and production
capacity to maximize profit or minimize cost.

\medskip

\textbf{Transportation.}
Shipping plans minimize transportation cost while satisfying supply and
demand constraints.

\medskip

\textbf{Supply Chains.}
Linear models allocate products among warehouses, distribution centers,
and customers.

\medskip

\textbf{Blending.}
Raw materials are combined to satisfy quality specifications at minimum
cost.

\medskip

\textbf{Agriculture.}
Land, water, labor, feed, and crop decisions can be optimized using LP
models.

\medskip

\textbf{Energy Systems.}
Linear programs model generation scheduling, fuel allocation, and
simplified power-flow decisions.

\medskip

\textbf{Finance.}
Some portfolio, cash-flow, and arbitrage models can be formulated as
linear programs.

\medskip

\textbf{Public Planning.}
Budget allocation, staffing, emergency logistics, and public-resource
planning frequently use linear optimization.

\end{applicationbox}

%%%%%%%%%%%%%%%%%%%%%%%%%%%%%%%%%%%%%%%%%%%%%%%%%%%%%%%%%%%%
\subsection{Exercises}
\label{subsec:linear-programming-exercises}
%%%%%%%%%%%%%%%%%%%%%%%%%%%%%%%%%%%%%%%%%%%%%%%%%%%%%%%%%%%%

\begin{exercise}[1]
Define a linear programming problem.
\end{exercise}

\begin{exercise}[1]
Define a polyhedron.
\end{exercise}

\begin{exercise}[1]
Define an extreme point.
\end{exercise}

\begin{exercise}[1]
Define a slack variable.
\end{exercise}

\begin{exercise}[1]
Define a surplus variable.
\end{exercise}

\begin{exercise}[1]
Define a basic feasible solution.
\end{exercise}

\begin{exercise}[1]
Define degeneracy.
\end{exercise}

\begin{exercise}[1]
Define the dual of a linear program conceptually.
\end{exercise}

\begin{exercise}[1]
State weak duality.
\end{exercise}

\begin{exercise}[1]
State strong duality.
\end{exercise}

\begin{exercise}[1]
State complementary slackness.
\end{exercise}

\begin{exercise}[2]
Convert

\[
2x_1+x_2\leq8
\]

to equality form using a slack variable.
\end{exercise}

\begin{exercise}[2]
Convert

\[
3x_1+4x_2\geq12
\]

to equality form using a surplus variable.
\end{exercise}

\begin{exercise}[2]
Represent an unrestricted variable \(x\) using nonnegative variables.
\end{exercise}

\begin{exercise}[2]
Determine whether the feasible region

\[
x+y\leq5,
\qquad
x,y\geq0
\]

is bounded.
\end{exercise}

\begin{exercise}[2]
Find all vertices of

\[
x+y\leq4,
\]

\[
x\leq3,
\]

\[
y\leq2,
\]

\[
x,y\geq0.
\]
\end{exercise}

\begin{exercise}[2]
Maximize

\[
z=2x+3y
\]

over the feasible region in the previous problem.
\end{exercise}

\begin{exercise}[2]
Give an example of an infeasible LP.
\end{exercise}

\begin{exercise}[2]
Give an example of an unbounded LP.
\end{exercise}

\begin{exercise}[2]
Give an example of an LP with multiple optimal solutions.
\end{exercise}

\begin{exercise}[3]
Prove that the feasible set of a linear program is convex.
\end{exercise}

\begin{exercise}[3]
Explain why a finite attained LP optimum can be found at an extreme point.
\end{exercise}

\begin{exercise}[3]
For

\[
A=
\begin{pmatrix}
1&1&1&0\\
2&1&0&1
\end{pmatrix},
\]

identify possible basis matrices.
\end{exercise}

\begin{exercise}[3]
Compute a basic solution for one basis chosen in the previous problem.
\end{exercise}

\begin{exercise}[3]
Determine whether that basic solution is feasible.
\end{exercise}

\begin{exercise}[3]
Explain how degeneracy occurs at an LP vertex.
\end{exercise}

\begin{exercise}[3]
Perform one simplex pivot on a small tableau.
\end{exercise}

\begin{exercise}[3]
Explain the entering-variable criterion for a maximization simplex
tableau.
\end{exercise}

\begin{exercise}[3]
Explain the ratio test geometrically.
\end{exercise}

\begin{exercise}[3]
Explain how the simplex method detects unboundedness.
\end{exercise}

\begin{exercise}[3]
Construct a Phase I auxiliary problem for an equality-constrained LP.
\end{exercise}

\begin{exercise}[3]
Explain why a positive Phase I optimum proves infeasibility.
\end{exercise}

\begin{exercise}[3]
Construct the dual of

\[
\begin{aligned}
\max\quad&
4x_1+5x_2
\\
\text{subject to}\quad&
2x_1+x_2\leq10,
\\
&
x_1+3x_2\leq12,
\\
&
x_1,x_2\geq0.
\end{aligned}
\]
\end{exercise}

\begin{exercise}[3]
Prove weak duality for the standard primal-dual pair.
\end{exercise}

\begin{exercise}[3]
Use complementary slackness to recover a dual solution from a known
primal optimum.
\end{exercise}

\begin{exercise}[3]
Use equal primal and dual objective values to certify optimality.
\end{exercise}

\begin{exercise}[3]
Interpret a positive dual variable as a shadow price.
\end{exercise}

\begin{exercise}[3]
Explain why an inactive primal resource constraint has zero dual price
under complementary slackness.
\end{exercise}

\begin{exercise}[3]
Compute the reduced costs for a specified basis.
\end{exercise}

\begin{exercise}[3]
Identify an alternate optimum using a zero nonbasic reduced cost.
\end{exercise}

\begin{exercise}[3]
Formulate a two-origin, three-destination transportation problem.
\end{exercise}

\begin{exercise}[3]
Formulate a diet problem with three foods and two nutritional
requirements.
\end{exercise}

\begin{exercise}[3]
Formulate a blending problem with quality constraints.
\end{exercise}

\begin{exercise}[3]
Formulate a production-planning problem with labor and material
constraints.
\end{exercise}

\begin{exercise}[4]
Prove the equivalence between basic feasible solutions and extreme points
under standard full-row-rank assumptions.
\end{exercise}

\begin{exercise}[4]
Study the simplex method algebraically using basis matrices.
\end{exercise}

\begin{exercise}[4]
Derive the reduced-cost formula from the simplex dictionary.
\end{exercise}

\begin{exercise}[4]
Study degeneracy and cycling.
\end{exercise}

\begin{exercise}[4]
Prove Bland's anti-cycling rule.
\end{exercise}

\begin{exercise}[4]
Derive strong LP duality.
\end{exercise}

\begin{exercise}[4]
Derive complementary slackness from strong duality.
\end{exercise}

\begin{exercise}[4]
Study Farkas' lemma and its equivalence to LP duality.
\end{exercise}

\begin{exercise}[4]
Construct infeasibility certificates using dual multipliers.
\end{exercise}

\begin{exercise}[4]
Study sensitivity analysis for right-hand-side perturbations.
\end{exercise}

\begin{exercise}[4]
Study allowable objective-coefficient ranges for a fixed optimal basis.
\end{exercise}

\begin{exercise}[4]
Study total unimodularity.
\end{exercise}

\begin{exercise}[4]
Prove integrality of network-flow polyhedra under suitable assumptions.
\end{exercise}

\begin{exercise}[4]
Study the transportation simplex method.
\end{exercise}

\begin{exercise}[4]
Study Dantzig--Wolfe decomposition.
\end{exercise}

\begin{exercise}[4]
Study column generation.
\end{exercise}

\begin{exercise}[4]
Derive primal-dual interior-point equations for LP.
\end{exercise}

\begin{exercise}[4]
Study the logarithmic-barrier central path.
\end{exercise}

\begin{exercise}[4]
Compare theoretical complexity of simplex and interior-point methods.
\end{exercise}

\begin{exercise}[4]
Investigate numerical scaling and presolve in large sparse LPs.
\end{exercise}

%%%%%%%%%%%%%%%%%%%%%%%%%%%%%%%%%%%%%%%%%%%%%%%%%%%%%%%%%%%%
\subsection{Conceptual Summary}
\label{subsec:linear-programming-summary}
%%%%%%%%%%%%%%%%%%%%%%%%%%%%%%%%%%%%%%%%%%%%%%%%%%%%%%%%%%%%

A standard linear program is

\[
\boxed{
\begin{aligned}
\max_{\mathbf{x}}
\quad&
\mathbf{c}^T\mathbf{x}
\\
\text{subject to}
\quad&
A\mathbf{x}\leq\mathbf{b},
\\
&
\mathbf{x}\geq0.
\end{aligned}
}
\]

Its feasible region is a convex polyhedron.

When a finite optimum is attained, an optimal extreme point exists.

Slack variables convert inequalities into equality form.

A basic feasible solution corresponds to a basis of the constraint matrix
and represents a vertex of the feasible polyhedron under standard
assumptions.

The simplex method moves between basic feasible solutions.

For the standard primal

\[
\boxed{
\max
\{
c^Tx:
Ax\leq b,
x\geq0
\},
}
\]

the dual is

\[
\boxed{
\min
\{
b^Ty:
A^Ty\geq c,
y\geq0
\}.
}
\]

Weak duality gives

\[
\boxed{
c^Tx
\leq
b^Ty.
}
\]

Strong duality gives equality at optimality:

\[
\boxed{
c^Tx^*
=
b^Ty^*.
}
\]

Complementary slackness links primal activity with dual variables.

Linear programming therefore combines

\[
\boxed{
\text{linear algebra}
+
\text{convex geometry}
+
\text{duality}
+
\text{algorithms}
+
\text{decision modeling}.
}
\]

%%%%%%%%%%%%%%%%%%%%%%%%%%%%%%%%%%%%%%%%%%%%%%%%%%%%%%%%%%%%
\subsection{Section Connections}
\label{subsec:linear-programming-connections}
%%%%%%%%%%%%%%%%%%%%%%%%%%%%%%%%%%%%%%%%%%%%%%%%%%%%%%%%%%%%

\chapterconnection{
Linear Programming specializes the general optimization framework of
Section~\ref{sec:optimization-fundamentals} to linear objectives and
polyhedral feasible regions. Extreme points lead to the simplex method,
while duality gives optimality certificates, shadow prices, and
sensitivity information. These ideas become foundational for integer,
combinatorial, network, scheduling, and stochastic optimization. The next
section, Nonlinear Programming, removes the linearity restriction and
develops general smooth constrained optimization, second-order
conditions, nonlinear KKT systems, penalty methods, sequential quadratic
programming, and related algorithms.
}

%%%%%%%%%%%%%%%%%%%%%%%%%%%%%%%%%%%%%%%%%%%%%%%%%%%%%%%%%%%%
% SECTION SOURCE TRACKING
%%%%%%%%%%%%%%%%%%%%%%%%%%%%%%%%%%%%%%%%%%%%%%%%%%%%%%%%%%%%

%------------------------------------------------------------
% 10.2 Linear Programming
%------------------------------------------------------------
%
% Source basis:
%   - Standard linear programming theory.
%   - Standard simplex and duality theory.
%   - Standard operations-research modeling.
%
% Used for:
%   - LP definitions and matrix form;
%   - feasible regions and polyhedra;
%   - convexity;
%   - extreme points;
%   - graphical LP methods;
%   - boundedness and unboundedness;
%   - infeasibility;
%   - alternate optima;
%   - standard form;
%   - slack, surplus, and unrestricted variables;
%   - basic solutions;
%   - basic feasible solutions;
%   - basis matrices;
%   - degeneracy;
%   - simplex method;
%   - reduced costs;
%   - entering and leaving variables;
%   - ratio test;
%   - pivots;
%   - simplex optimality;
%   - simplex unboundedness;
%   - cycling and Bland's rule;
%   - artificial variables;
%   - Phase I and Phase II;
%   - Big-M preview;
%   - LP duality;
%   - weak and strong duality;
%   - duality gaps;
%   - complementary slackness;
%   - shadow prices;
%   - sensitivity analysis;
%   - Farkas' lemma preview;
%   - optimality and infeasibility certificates;
%   - transportation, blending, diet, and production models;
%   - network-flow preview;
%   - total unimodularity preview;
%   - interior-point methods;
%   - central path;
%   - computational complexity;
%   - presolve and scaling;
%   - applications and exercises.
%
% Notes:
%   - Nonlinear Programming is developed next in Section 10.3.
%   - Convex Optimization follows in Section 10.4.
%   - Integer, combinatorial, and network optimization build directly
%     on linear-programming ideas later in Chapter 10.
%
%%%%%%%%%%%%%%%%%%%%%%%%%%%%%%%%%%%%%%%%%%%%%%%%%%%%%%%%%%%%